\section{Proofs}
\label{app:proofs}

\subsection{Proofs for Section~\ref{sec:identification}}

\paragraph{Sufficient conditions for outcome-relevant completeness.}
\begin{proof}[Proof of Proposition~\ref{appdetail:prop:orc-sufficient}]
Fix $z$ and $Q,Q'\in\mathcal Q_{j,z}$ with $\mathcal K_jQ=\mathcal K_jQ'$. Write $\pi_Q$ for the \textcolor{red}{\sout{law}} \textcolor{blue}{distribution} of $(X,U)$ under $Q$, $\pi_{Q,X}$ for its $X$-marginal, and $\pi_Q(\cdot\mid x)$ for its conditional \textcolor{red}{\sout{law}} \textcolor{blue}{distribution} of $U$ given $X=x$. By condition (i), $K_j(\cdot\mid v,u)=K_j(\cdot\mid x,u)$, so the \textcolor{red}{\sout{law}} \textcolor{blue}{distribution} $\mathcal K_jQ$ of $(X,W_j)$ has $X$-marginal $\pi_{Q,X}$ and conditional \textcolor{red}{\sout{law}} \textcolor{blue}{distribution} $\int K_j(\cdot\mid x,u)\,\pi_Q(du\mid x)$ given $X=x$. The equality $\mathcal K_jQ=\mathcal K_jQ'$ therefore gives $\pi_{Q,X}=\pi_{Q',X}$ and, for $\pi_{Q,X}$-almost every $x$, $\int K_j(\cdot\mid x,u)\,\pi_Q(du\mid x)=\int K_j(\cdot\mid x,u)\,\pi_{Q'}(du\mid x)$. Condition (ii) yields $\pi_Q(\cdot\mid x)=\pi_{Q'}(\cdot\mid x)$ for almost every $x$, hence $\pi_Q=\pi_{Q'}$. By condition (i), $\mu_{a,j,z}(v,u)=\bar\mu_{a,z}(x,u)$ for a measurable function $\bar\mu_{a,z}$, so $\mathbb E_Q\{\mu_{a,j,z}(V_j,U)\}=\int\bar\mu_{a,z}\,d\pi_Q=\int\bar\mu_{a,z}\,d\pi_{Q'}=\mathbb E_{Q'}\{\mu_{a,j,z}(V_j,U)\}$, which is Assumption~\ref{ass:outcome-relevant-completeness}.

For case (a), the \textcolor{red}{\sout{law}} \textcolor{blue}{distribution} of $W_j$ given $X=x$ is the convolution of the pushforward of $\pi$ under $g(x,\cdot)$ with the \textcolor{red}{\sout{law}} \textcolor{blue}{distribution} of $\epsilon_j$. Its characteristic function is the product of the characteristic function of the pushforward and that of $\epsilon_j$, and the latter vanishes nowhere by assumption. Two \textcolor{red}{\sout{laws}} \textcolor{blue}{distributions} $\pi$ and $\pi'$ with equal convolutions therefore have pushforwards with equal characteristic functions, hence equal pushforwards by the uniqueness theorem for characteristic functions, and injectivity of $g(x,\cdot)$ gives $\pi=\pi'$. Thus (ii) holds.

For case (b), the probability generating function of $W_j$ given $X=x$ under a \textcolor{red}{\sout{law}} \textcolor{blue}{distribution} $\pi$ of $U$ is $\mathbb E(\prod_\ell s_\ell^{W_{j,\ell}})=\int\exp\{-\sum_\ell\lambda_\ell(x,u)(1-s_\ell)\}\,\pi(du)$ for $s\in[0,1]^{d}$, which is the Laplace transform of the pushforward of $\pi$ under $\lambda(x,\cdot)$ evaluated on $[0,1]^{d}$. A finite measure on $[0,\infty)^{d}$ is determined by its Laplace transform on a set with nonempty interior, so two \textcolor{red}{\sout{laws}} \textcolor{blue}{distributions} with equal mixtures have equal pushforwards, and injectivity of $\lambda(x,\cdot)$ gives $\pi=\pi'$. In the scalar case this is the identifiability of Poisson mixtures of \citet{teicher1961identifiability}. Thus (ii) holds.

\end{proof}

\paragraph{Exact identification.}
\begin{proof}[Proof of Theorem~\ref{thm:exact-proxy-identification}]
Because $Z_j=\phi_j(V_j)$ is a measurable function of $V_j$, conditioning on $(U,V_j,Z_j)$ is equivalent to conditioning on $(U,V_j)$. The two clauses of Assumption~\ref{ass:joint-latent-exchangeability} therefore remain valid after additionally conditioning on $Z_j$. Thus, for each $a\in\{0,1\}$,
\begin{align}
\mathbb E\{Y(a)\mid A,V_j,U,Z_j\}
&=
\mathbb E\{Y(a)\mid V_j,U,Z_j\},
\label{eq:conditioning-safe-outcome-app}\\
 \textcolor{red}{\mathcal L}\textcolor{blue}{P}(W_j\mid A,V_j,U,Z_j)
&=
K_j(\cdot\mid V_j,U).
\nonumber
\end{align}
Integrating the second identity over $Q_{b,j,z}$ and using the definition of $\mathcal K_j$ gives
\( \textcolor{red}{\mathcal L}\textcolor{blue}{P}(X,W_j\mid A=b,Z_j=z) = \mathcal K_jQ_{b,j,z}.\)
Exact balance in Eq.~\eqref{eq:exact-proxy-balance} therefore implies
\(\mathcal K_jQ_{1,j,z}=\mathcal K_jQ_{0,j,z}\)
for $P_{Z_j}$-almost every $z$. Representation overlap gives $Q_{b,j,z}\ll \textcolor{red}{\mathcal L}\textcolor{blue}{P}(V_j,U\mid Z_j=z)$ for $b\in\{0,1\}$. Assumption~\ref{ass:consistency}, Eq.~\eqref{eq:conditioning-safe-outcome-app}, and Jensen's inequality give the required integrability, so $Q_{b,j,z}\in\mathcal Q_{j,z}$ almost everywhere. Assumption~\ref{ass:outcome-relevant-completeness} therefore gives, for each $a\in\{0,1\}$,
\begin{align}
\mathbb E_{Q_{1,j,z}}\{\mu_{a,j,z}(V_j,U)\}
=
\mathbb E_{Q_{0,j,z}}\{\mu_{a,j,z}(V_j,U)\}.
\label{eq:orc-equal-latent-means}
\end{align}
By Eq.~\eqref{eq:conditioning-safe-outcome-app} and iterated expectation, for each $b\in\{0,1\}$,
\(\mathbb E\{Y(a)\mid A=b,Z_j=z\} = \mathbb E_{Q_{b,j,z}}\{\mu_{a,j,z}(V_j,U)\}.\)
Eq.~\eqref{eq:orc-equal-latent-means} shows that the conditional mean of $Y(a)$ given $(A,Z_j)$ does not depend on $A$. Consistency then gives
\(\mathbb E(Y\mid A=a,Z_j=z) = \mathbb E\{Y(a)\mid A=a,Z_j=z\} = \mathbb E\{Y(a)\mid Z_j=z\}.\)
Integrating over $P_{Z_j}$ yields Eq.~\eqref{eq:exact-identification}.\qedhere
\end{proof}

\paragraph{Non-nesting of canonical proximal and PROBE sufficient conditions.}
\begin{proof}[Proof of Proposition~\ref{prop:non-nesting-proximal-probe}]
We give one finite-state \textcolor{red}{\sout{law}} \textcolor{blue}{distribution} for each direction. In both constructions, $X$ is constant, $Y=Y(A)$, and every treatment draw is generated by an exogenous randomizer independent of all proxy noises and outcome noises.

First, we construct a \textcolor{red}{\sout{law}} \textcolor{blue}{distribution} satisfying the canonical proximal conditions but not the PROBE conditions for the designated proxy split. Let
\(U\sim\operatorname{Bernoulli}(1/2), P(A=1\mid U=u)=\frac14+\frac12u,\)
and let
\(Z^{\mathrm{tr}}=U\mathbin\oplus E_Z, \qquad Z^{\mathrm{out}}=W_j=U\mathbin\oplus E_W,\)
where $E_Z$ and $E_W$ are mutually independent $\operatorname{Bernoulli}(1/4)$ variables, independent of the treatment randomizer. Set $Y(a)=a+U$ and $V_j=Z^{\mathrm{tr}}$. Consistency, integrability, and strict latent overlap hold. Independence of the exogenous noises gives the three canonical proxy-validity conditions.

The function
\(h(z^{\mathrm{out}},a)=a+2z^{\mathrm{out}}-\frac12\)
is an outcome bridge because $\mathbb E(Z^{\mathrm{out}}\mid U=u)=1/4+u/2$, so $\mathbb E\{h(Z^{\mathrm{out}},a)\mid U=u\}=a+u$. Iterated expectation gives the required observed bridge equation. The binary channel from $U$ to $Z^{\mathrm{tr}}$ has matrix
\(\begin{pmatrix} 3/4&1/4\\ 1/4&3/4 \end{pmatrix}\)
and determinant $1/2$. Conditional on either treatment arm, both latent states have positive probability, and the corresponding unnormalized conditional-expectation matrix has determinant $3/32$. It therefore has full column rank, which proves the canonical treatment-proxy completeness condition.

Now consider any measurable representation $Z_j=\phi_j(V_j)$. The Markov relation $U\to V_j\to Z_j$ holds. Conditional on $(U,Z_j)$, treatment and $W_j$ are independent, and
\(\mathbb E(A\mid U,Z_j) = \mathbb E(W_j\mid U,Z_j) = \frac14+\frac12U.\)
The conditional covariance decomposition gives
\(\operatorname{Cov}(A,W_j\mid Z_j) = \frac14\operatorname{Var}(U\mid Z_j).\)
Because $P(U=1\mid V_j=0)=1/4$ and $P(U=1\mid V_j=1)=3/4$,
\(P(U=1\mid Z_j) = \mathbb E\{P(U=1\mid V_j)\mid Z_j\} \in[1/4,3/4]\)
almost surely. Hence $\operatorname{Var}(U\mid Z_j)\ge3/16$ and
\(\operatorname{Cov}(A,W_j\mid Z_j)\ge\frac{3}{64}>0.\)
Exact balance therefore fails for every such representation. This proves the first direction.

Second, we construct a \textcolor{red}{\sout{law}} \textcolor{blue}{distribution} satisfying the PROBE conditions for the designated proxy split but not the canonical proximal conditions. Let $U=(U_1,U_2)$, where $U_1$ and $U_2$ are independent $\operatorname{Bernoulli}(1/2)$ variables. Let $Z^{\mathrm{tr}}$ and $Z^{\mathrm{out}}$ be separate observed coordinates that agree almost surely with $U_1$, and set
\(V_j=Z^{\mathrm{tr}}=U_1, \qquad W_j=Z^{\mathrm{out}}=U_1, \qquad Z_j=U_1.\)
Define
\(P(A=1\mid U_1,U_2) = \frac3{20}+\frac15U_1+\frac25U_2,\)
and let $\xi$ be an independent mean-zero Rademacher variable, also independent of the treatment randomizer. Set
\(Y(a)=a+U_1+(1+U_2)\xi.\)
The four latent treatment probabilities are $3/20$, $11/20$, $7/20$, and $15/20$, so strict latent overlap holds. Latent exchangeability follows from the independent outcome noise, and held-out treatment independence holds because $W_j$ is deterministic given $U$. Moreover,
\(P(A=1\mid Z_j=0)=\frac7{20}, \qquad P(A=1\mid Z_j=1)=\frac{11}{20},\)
so representation overlap holds. Because $W_j=Z_j$ and $X$ is constant, exact balance also holds.

It remains to verify the all-$Q$ outcome-relevant completeness condition. Within a stratum $Z_j=z$, every admissible \textcolor{red}{\sout{law}} \textcolor{blue}{distribution} $Q\in\mathcal Q_{j,z}$ is supported on $V_j=U_1=z$, while $U_2$ may vary. On this support,
\(\mu_{a,j,z}(v,u)=a+z.\)
Thus, for any admissible $Q,Q'$ and either potential treatment level,
\(\mathcal K_jQ=\mathcal K_jQ' \quad\Longrightarrow\quad \mathbb E_Q\{\mu_{a,j,z}(V_j,U)\} =a+z =\mathbb E_{Q'}\{\mu_{a,j,z}(V_j,U)\}.\)
This proves all PROBE conditions for the designated proxy split.

The outcome bridge $h(z^{\mathrm{out}},a)=a+z^{\mathrm{out}}$ exists because
\(\mathbb E(Y\mid Z^{\mathrm{tr}}=z,A=a) =a+z =\mathbb E\{h(Z^{\mathrm{out}},a)\mid Z^{\mathrm{tr}}=z,A=a\}.\)
Canonical treatment-proxy completeness nevertheless fails. In the treated arm,
\(P(U_2=1\mid U_1=0,A=1)=\frac{11}{14}, \qquad P(U_2=1\mid U_1=1,A=1)=\frac{15}{22}.\)
Define the square-integrable, nonzero function
\(g_1(U) = U_2- \begin{cases} 11/14,&U_1=0,\\ 15/22,&U_1=1. \end{cases}\)
Then
\(\mathbb E\{g_1(U)\mid Z^{\mathrm{tr}},A=1\} = \mathbb E\{g_1(U)\mid U_1,A=1\} =0\)
almost surely, although $g_1(U)$ is not almost surely zero. Hence the canonical completeness condition fails, proving the second direction.

For the final claim, the second construction satisfies Assumption~\ref{ass:outcome-relevant-completeness} but not canonical treatment-proxy completeness. In the first construction, replace $E_W$ by an independent $\operatorname{Bernoulli}(1/2)$ variable. Canonical treatment-proxy completeness is unchanged because it involves only $(U,Z^{\mathrm{tr}},A)$. Now $W_j$ is independent of $U$, so $\mathcal K_jQ$ is the same \textcolor{red}{\sout{law}} \textcolor{blue}{distribution} for every $Q$. For every representation $\phi_j$ and almost every $z$, the bound $P(U=1\mid Z_j)\in[1/4,3/4]$ shown above gives both values of $U$ positive probability given $Z_j=z$, so $\mathbb E_Q\mu_{a,j,z}=a+Q(U=1)$ takes more than one value over $\mathcal Q_{j,z}$ and Assumption~\ref{ass:outcome-relevant-completeness} fails.
\end{proof}

\paragraph{Exhaustive plurality identification.}
\begin{proof}[Proof of Theorem~\ref{thm:unknown-valid-block-identification}]
Under condition~\eqref{eq:unique-target-plurality}, the number of proxy splits in $\mathcal B$ whose candidate value equals $\tau$ is strictly larger than the count for every non-target value. Therefore, $\tau$ is the unique value maximizing the population count:
\(\{\tau\} = \underset{c\in\mathbb R}{\operatorname{arg\,max}} \bigl|\{S\in\mathcal B:\theta_S=c\}\bigr|.\)
If no non-target value occurs, the convention in \textcolor{red}{\sout{Thm.}~\ref{thm:unknown-valid-block-identification}}\textcolor{blue}{Definition~\ref{def:population-plurality}} compares the nonempty target cluster with zero and gives the same conclusion.
\end{proof}

\paragraph{Target-valued majority.}
\begin{proof}[Proof of Corollary~\ref{cor:target-majority-median}]
Define
\(N_\tau \triangleq \bigl|\{S\in\mathcal B:\theta_S=\tau\}\bigr|,\); \(N_c \triangleq \bigl|\{S\in\mathcal B:\theta_S=c\}\bigr|, \qquad c\ne\tau.\)
Condition~\eqref{eq:target-valued-majority} gives $N_\tau>|\mathcal B|/2$. For every $c\ne\tau$,
\(N_c \le |\mathcal B|-N_\tau < N_\tau.\)
Therefore, $\tau$ is the unique candidate value of greatest multiplicity.

It remains to establish the median statement. Let $n=|\mathcal B|$. If $n=2k+1$ is odd, then $N_\tau\ge k+1$. At most $k$ candidate values are strictly below $\tau$, and at most $k$ are strictly above $\tau$, so the $(k+1)$st order statistic equals $\tau$. If $n=2k$ is even, then $N_\tau\ge k+1$ and $n-N_\tau\le k-1$. Hence both the $k$th and the $(k+1)$st order statistics equal $\tau$. Under the conventional definition of the even-sample median as the average of these two middle order statistics, the median equals $\tau$. Thus,
\(\operatorname{median}\bigl((\theta_S)_{S\in\mathcal B}\bigr)=\tau.\)
The stronger condition that more than half of the proxy splits in $\mathcal B$ satisfy the grouped assumptions is sufficient because each such proxy split has candidate value $\tau$, but it is not necessary.
\end{proof}

\paragraph{Zero discrepancy and exact balance.}
\begin{proof}[Proof of Proposition~\ref{prop:zero-discrepancy}]
By Definition~\ref{def:held-out-proxy-discrepancy}, $D_{\mathrm{res},j}(\phi)=0$ if and only if
\(q_{j,\phi}(T_j,Z_j^\phi)=e_{j,\phi}(Z_j^\phi) \quad\text{almost surely}.\)
Because $A$ is binary, this equality is equivalent to $T_j\perp\!\!\!\perp A\mid Z_j^\phi$.
\end{proof}

\paragraph{Supporting bias decomposition.}
Fix a representation $\phi$ and abbreviate $e(z)=e_{j,\phi}(z)$. For each $a\in\{0,1\}$ and $P_{Z_j^\phi}$-almost every $z$, define
\(\delta_{a,j}^{\phi}(z) \triangleq \mathbb E_{Q_{1,j,z}^{\phi}}\{\mu_{a,j,z}^{\phi}(V_j,U)\} - \mathbb E_{Q_{0,j,z}^{\phi}}\{\mu_{a,j,z}^{\phi}(V_j,U)\}.\)
Under Assumptions~\ref{ass:consistency} and \ref{ass:joint-latent-exchangeability}, and under $0<e(Z_j^\phi)<1$ almost surely, the conditioning argument used in Eq.~\eqref{eq:conditioning-safe-outcome-app} gives
\(\mathbb E(Y\mid A=a,Z_j^\phi=z) = \mathbb E_{Q_{a,j,z}^{\phi}}\{\mu_{a,j,z}^{\phi}(V_j,U)\}.\)
Conditional expectation over $A$ gives
\(\mathbb E\{Y(a)\mid Z_j^\phi=z\} ={} e(z)\mathbb E_{Q_{1,j,z}^{\phi}}\{\mu_{a,j,z}^{\phi}\}+\{1-e(z)\}\mathbb E_{Q_{0,j,z}^{\phi}}\{\mu_{a,j,z}^{\phi}\}.\)
Subtracting the two conditional expressions and collecting terms gives
\begin{align}
&\mathbb E(Y\mid A=1,Z_j^\phi=z)
-\mathbb E(Y\mid A=0,Z_j^\phi=z)
\nonumber\\
&\quad-
\left[
\mathbb E\{Y(1)\mid Z_j^\phi=z\}
-
\mathbb E\{Y(0)\mid Z_j^\phi=z\}
\right]
\nonumber\\
&\qquad=
\{1-e(z)\}\delta_{1,j}^{\phi}(z)
+e(z)\delta_{0,j}^{\phi}(z).
\label{eq:pointwise-bias-decomposition}
\end{align}

\paragraph{Approximate PROBE bias bound.}
\begin{proof}[Proof of Theorem~\ref{thm:approximate-proxy-bias-bound}]
Fix $z$ in the almost-sure set on which overlap and the structural conditions hold. Bayes' rule gives
\(\frac{d \textcolor{red}{\mathcal L}\textcolor{blue}{P}(T_j\mid A=1,Z_j^\phi=z)}{dP_{j,\phi,z}}(t) =\frac{q_{j,\phi}(t,z)}{e(z)},\); \(\frac{d \textcolor{red}{\mathcal L}\textcolor{blue}{P}(T_j\mid A=0,Z_j^\phi=z)}{dP_{j,\phi,z}}(t) =\frac{1-q_{j,\phi}(t,z)}{1-e(z)}.\)
Subtracting gives the density identity
\(\frac{d\{ \textcolor{red}{\mathcal L}\textcolor{blue}{P}(T_j\mid A=1,Z_j^\phi=z) - \textcolor{red}{\mathcal L}\textcolor{blue}{P}(T_j\mid A=0,Z_j^\phi=z)\}} {dP_{j,\phi,z}}(t) = \frac{q_{j,\phi}(t,z)-e(z)} {e(z)\{1-e(z)\}}.\)
Held-out treatment independence makes the conditional \textcolor{red}{\sout{law}} \textcolor{blue}{distribution} of $W_j$ given $(U,V_j)$ treatment-invariant, so $\mathcal K_jQ_{b,j,z}^{\phi}= \textcolor{red}{\mathcal L}\textcolor{blue}{P}(T_j\mid A=b,Z_j^\phi=z)$ for $b\in\{0,1\}$. As in the proof of Thm.~\ref{thm:exact-proxy-identification}, overlap and Assumption~\ref{ass:consistency} give $Q^{\phi}_{b,j,z}\in\mathcal Q^{\phi}_{j,z}$ for $b\in\{0,1\}$. Assumption~\ref{ass:uniform-stability} therefore implies
\begin{align}
|\delta_{a,j}^{\phi}(z)|
&\le
\frac{\Gamma_j(\phi)}{e(z)\{1-e(z)\}}
\left[
\mathbb E\!\left(
\left.
\{q_{j,\phi}(T_j,z)-e(z)\}^2
\right|Z_j^\phi=z
\right)
\right]^{1/2}.
\label{eq:stratum-bias-stability-bound}
\end{align}
Apply the triangle inequality to Eq.~\eqref{eq:pointwise-bias-decomposition}, use Eq.~\eqref{eq:stratum-bias-stability-bound} for both potential outcomes, and integrate over $Z_j^\phi$. This gives
\(\left| \mathbb E\!\left[ \mathbb E(Y\mid A=1,Z_j^\phi) - \mathbb E(Y\mid A=0,Z_j^\phi) \right] -\tau \right| \le \mathbb E\!\left[ \frac{\Gamma_j(\phi)} {e_{j,\phi}(Z_j^\phi)\{1-e_{j,\phi}(Z_j^\phi)\}} \left\{ \mathbb E\!\left[ \left. \{q_{j,\phi}(T_j,Z_j^\phi)-e_{j,\phi}(Z_j^\phi)\}^2 \right|Z_j^\phi \right] \right\}^{1/2} \right] \le \frac{\Gamma_j(\phi)}{\eta(1-\eta)}D_{\mathrm{res},j}(\phi).\)
Because $\Gamma_j(\phi)$ does not depend on $z$, the overlap bound $\eta\le e_{j,\phi}(z)\le1-\eta$ gives the last inequality, together with $\mathbb E\{s(Z_j^\phi)\}\le[\mathbb E\{s(Z_j^\phi)^2\}]^{1/2}=D_{\mathrm{res},j}(\phi)$ by Cauchy--Schwarz, where $s(z)$ denotes the conditional root-mean-square factor in the display, and proves Eq.~\eqref{eq:uniform-approximate-proxy-bias-bound}.
\end{proof}

\paragraph{Population sensitivity region.}
\begin{proof}[Proof of Proposition~\ref{prop:population-sensitivity-region}]
Thm.~\ref{thm:approximate-proxy-bias-bound} gives $|\tau_\phi-\tau|\le\Gamma_j(\phi)D_{\mathrm{res},j}(\phi)/\{\eta(1-\eta)\}$, and the displayed interval is the set of values of $\tau$ compatible with this inequality.
\end{proof}

\subsection{Proofs for Section~\ref{sec:estimation}}

\paragraph{Uniform discrepancy deviation.}
\begin{proof}[Proof of Proposition~\ref{appdetail:prop:uniform-discrepancy-deviation}]
For clarity, the joint encoder--critic loss classes used in the proposition are
\(\mathcal L_{0,j,n_{\mathrm D}} \triangleq \left\{ o\mapsto\{a-h(\phi(v_j))\}^2: \phi\in\mathcal F_{j,n_{\mathrm D}},\ h\in\mathcal H_{0,n_{\mathrm D}} \right\},\); \(\mathcal L_{1,j,n_{\mathrm D}} \triangleq \left\{ o\mapsto\{a-h(t_j,\phi(v_j))\}^2: \phi\in\mathcal F_{j,n_{\mathrm D}},\ h\in\mathcal H_{1,n_{\mathrm D}} \right\},\)
where $o=(v_j,t_j,a,y)$. For a loss class $\mathcal L$, its expected Rademacher complexity is
\(\mathfrak R_{n_{\mathrm D}}(\mathcal L) \triangleq \mathbb E\!\left[ \sup_{\ell\in\mathcal L} \frac1{n_{\mathrm D}}\sum_{i=1}^{n_{\mathrm D}}\sigma_i\ell(O_i) \right],\)
where the $\sigma_i$ are independent Rademacher signs independent of the observations. The loss classes are pointwise measurable and separable, so these suprema are measurable.
The same argument may instead be stated with outer probability and outer expectation.

Every loss in either class lies in $[0,1]$. Standard symmetrization and bounded-difference concentration therefore imply, simultaneously for $b\in\{0,1\}$, that
\begin{align}
\sup_{\phi\in\mathcal F_{j,n_{\mathrm D}},\,h\in\mathcal H_{b,n_{\mathrm D}}}
|\widehat R_{b,\phi}(h)-R_{b,\phi}(h)|
\le
2\mathfrak R_{n_{\mathrm D}}(\mathcal L_{b,j,n_{\mathrm D}})
+\sqrt{\frac{\log(4/\delta_{\mathrm D})}{2n_{\mathrm D}}}
\label{eq:brier-risk-uniform-deviation-app}
\end{align}
with probability at least $1-\delta_{\mathrm D}$. Let $R_{b,\phi}^{\star}=\inf_{h\in\mathcal H_{b,n_{\mathrm D}}}R_{b,\phi}(h)$ and define $\widehat R_{b,\phi}^{\star}$ analogously. Eq.~\eqref{eq:brier-risk-uniform-deviation-app} controls $|\widehat R_{b,\phi}^{\star}-R_{b,\phi}^{\star}|$ by the same right-hand side. Eq.~\eqref{eq:approximate-critic-fits} also gives
\(0 \le \widehat R_{b,\phi}(\widehat h_{b,\phi}) -\widehat R_{b,\phi}^{\star} \le \varepsilon_{b,\phi}.\)

The conditional-mean property of squared loss yields the Brier identity
\(\mathbb E[\{A-e_{j,\phi}(Z_j^\phi)\}^2] -\mathbb E[\{A-q_{j,\phi}(T_j,Z_j^\phi)\}^2] = \mathbb E[\{q_{j,\phi}(T_j,Z_j^\phi)-e_{j,\phi}(Z_j^\phi)\}^2] =D_{\mathrm{res},j}^2(\phi).\)
It follows from Eq.~\eqref{eq:critic-approximation-error} that
\begin{align}
\left|
R_{0,\phi}^{\star}-R_{1,\phi}^{\star}
-D_{\mathrm{res},j}^2(\phi)
\right|
\le
a_{0,n_{\mathrm D}}(\phi)+a_{1,n_{\mathrm D}}(\phi)
\le a_{j,n_{\mathrm D}}.
\label{eq:population-critic-gap-app}
\end{align}
Combine Eqs.~\eqref{eq:brier-risk-uniform-deviation-app}--\eqref{eq:population-critic-gap-app}. The positive-part map is one-Lipschitz and $D_{\mathrm{res},j}^2(\phi)\ge0$, so taking the supremum over $\phi$ gives Eq.~\eqref{eq:uniform-discrepancy-deviation}.
\end{proof}

\paragraph{Representation-learning oracle inequality.}
\begin{proof}[Proof of Proposition~\ref{appdetail:thm:representation-oracle}]
On the event in Eq.~\eqref{eq:uniform-discrepancy-deviation}, set
\(\xi_{j,n_{\mathrm D}} \triangleq r_{j,n_{\mathrm D}}(\delta_{\mathrm D})+a_{j,n_{\mathrm D}}+\varepsilon_{\mathrm{crit},j}.\)
The deviation inequality at $\widehat\phi_j$, Eq.~\eqref{eq:approximate-erm}, and the deviation inequality at a near-minimizer of $D_{\mathrm{res},j}^2$ over $\mathcal F_{j,n_{\mathrm D}}$ give
\(D_{\mathrm{res},j}^2(\widehat\phi_j) \le \widetilde D_{j,n_{\mathrm D}}^2(\widehat\phi_j)+\xi_{j,n_{\mathrm D}} \le \inf_{\phi\in\mathcal F_{j,n_{\mathrm D}}} \widetilde D_{j,n_{\mathrm D}}^2(\phi) +\varepsilon_{\mathrm{enc},j}+\xi_{j,n_{\mathrm D}} \le \inf_{\phi\in\mathcal F_{j,n_{\mathrm D}}} D_{\mathrm{res},j}^2(\phi) +2\xi_{j,n_{\mathrm D}}+\varepsilon_{\mathrm{enc},j}\),
which proves Eq.~\eqref{eq:representation-oracle-discrepancy}.
\end{proof}

\paragraph{Causal error of the learned representation.}
\begin{proof}[Proof of Theorem~\ref{cor:representation-causal}]
Conditional on the training observations $\{O_i:i\in\mathcal I_{\mathrm D}\}$ and any algorithmic seed used in training, the selected encoder is fixed; the population expectations below use a fresh observation independent of these training inputs. Eq.~\eqref{eq:uniform-approximate-proxy-bias-bound} gives
\(|\tau_{\widehat Z_j}-\tau| \le \frac{\Gamma_j(\widehat\phi_j)}{\eta(1-\eta)} D_{\mathrm{res},j}(\widehat\phi_j).\)
Taking the square root of Eq.~\eqref{eq:representation-oracle-discrepancy} and substituting it into this display proves Eq.~\eqref{eq:representation-oracle-causal}.
\end{proof}

\paragraph{Exact signed nuisance remainder.}
\begin{proof}[Proof of Proposition~\ref{appdetail:prop:aipw-remainder}]
Fix $\widehat Z_j=z$. Taking the conditional expectation of the score and using Eq.~\eqref{eq:aipw-nuisances} gives
\(\mathbb E\{\mathtt{UIF}_j(O)\mid\widehat Z_j=z\} ={} \widehat m_{1,j}(z)-\widehat m_{0,j}(z)+ \frac{e_j(z)\{m_{1,j}(z)-\widehat m_{1,j}(z)\}} {\widehat e_j(z)}- \frac{\{1-e_j(z)\}\{m_{0,j}(z)-\widehat m_{0,j}(z)\}} {1-\widehat e_j(z)}.\)
Subtract $m_{1,j}(z)-m_{0,j}(z)$ and collect the two nuisance errors. Averaging the resulting identity over $\widehat Z_j$ proves Eq.~\eqref{eq:aipw-remainder}.
\end{proof}

\paragraph{Conditional finite-sample AIPW error.}
\begin{proof}[Proof of Theorem~\ref{thm:conditional-aipw-error}]
 Conditional on the two learning samples, the evaluation observations are independent and the score $\mathtt{UIF}_j(O)$ has variance $\sigma_j^2$, so Chebyshev's inequality bounds the difference between the evaluation average and its conditional expectation by $\sigma_j/\sqrt{n_{\mathrm E}\delta_{\mathrm E}}$, the first term in Eq.~\eqref{eq:aipw-error-radius}, with probability at least $1-\delta_{\mathrm E}$. Prop.~\ref{appdetail:prop:aipw-remainder}, Cauchy--Schwarz, and the propensity clipping bound control the absolute conditional remainder by
\(\frac{q_{e,j}(q_{0,j}+q_{1,j})}{\eta}.\)
The triangle inequality proves Eq.~\eqref{eq:conditional-aipw-error}.
For Eq.~\eqref{eq:end-to-end-sensitivity}, the triangle inequality gives $|\widehat\tau_j^{\mathrm{AIPW}}-\tau|\le|\widehat\tau_j^{\mathrm{AIPW}}-\tau_{\widehat Z_j}|+|\tau_{\widehat Z_j}-\tau|$. Eq.~\eqref{eq:conditional-aipw-error} bounds the first term on the same event, and Eq.~\eqref{eq:uniform-approximate-proxy-bias-bound} bounds the second term by $\Gamma_j(\widehat\phi_j)D_{\mathrm{res},j}(\widehat\phi_j)/\{\eta(1-\eta)\}$.
\end{proof}

\paragraph{Consistency.}
\begin{proof}[Proof of Corollary~\ref{cor:learned-aipw-consistency}]
 Condition (C1) makes the evaluation term $\sigma_{j,n}/\sqrt{n_{\mathrm E}\delta_{\mathrm E,n}}$ of $\varepsilon_{\mathrm{AIPW},j}$ vanish, condition (C2) makes its nuisance-product term vanish, and condition (C3) makes the stability-weighted residual-discrepancy term in Eq.~\eqref{eq:end-to-end-sensitivity} vanish. The failure probability $\delta_{\mathrm E,n}$ of Thm.~\ref{thm:conditional-aipw-error} also converges to zero, which proves convergence in probability.
\end{proof}

\paragraph[Cluster plurality under random subsampling.]{ Cluster plurality under random subsampling.}
\begin{proof}[ Proof of Lemma~\ref{lem:random-subsampling-plurality}]
Conditional on $\mathcal G$ and $R=r\ge1$, the retained proxy splits form a uniformly sampled $r$-element subset of $\widehat{\mathcal B}$ because the $m$ draws are uniform without replacement and $\widehat{\mathcal B}$ is fixed. For each competing cell $C_\ell$, define $H_\ell(S)\triangleq\mathbf 1\{S\in C_0\}-\mathbf 1\{S\in C_\ell\}$ for $S\in\widehat{\mathcal B}$. This variable lies in $[-1,1]$ and has population mean $(n_0-n_\ell)/N\ge\Delta_\Pi$. Let $\overline H_{\ell,r}$ be its retained-sample average. The event $R_\ell\ge R_0$ implies $\overline H_{\ell,r}\le0$. Hoeffding's inequality for sampling without replacement \citep{hoeffding1963probability} gives $P(\overline H_{\ell,r}\le0\mid\mathcal G,R=r)\le e^{-r\Delta_\Pi^2/2}$. A union bound over the $L$ competing cells yields $P(\max_{1\le\ell\le L}R_\ell\ge R_0\mid\mathcal G,R=r)\le L e^{-r\Delta_\Pi^2/2}$. Averaging over $R$ and counting $R=0$ as failure gives $P(\text{failure}\mid\mathcal G)\le P(R=0\mid\mathcal G)+L\,\mathbb E\{e^{-R\Delta_\Pi^2/2}\mid\mathcal G\}$. If $L=0$, every retained split belongs to $C_0$, so the conclusion holds whenever $R\ge1$.
\end{proof}

\paragraph[Error of Algorithm 1.]{ Error of Algorithm~\ref{alg:probe-search}.}
\begin{proof}[ Proof of Theorem~\ref{thm:probe-search-error}]
Work conditionally on $\mathcal G$. Let $E_1$ be the uniform-evaluation-error event and let $E_2$ be the event that $R\ge1$ and $R_0>\max_{1\le\ell\le L}R_\ell$. \textcolor{red}{\sout{Condition (iii) of Thm.~}\ref{thm:probe-search-error}\sout{ gives $P(E_1^{\mathrm c}\mid\mathcal G)\le\delta$,}} \textcolor{blue}{Eq.~\eqref{eq:conditional-aipw-error} applied to each $S\in\widehat{\mathcal B}$ at $\delta_{\mathrm E}=\delta/N$ gives $P(|\widehat\theta_S-\theta_S|>\varepsilon\mid\mathcal G)\le\delta/N$, because $\varepsilon$ is at least the radius of Eq.~\eqref{eq:aipw-error-radius} for $S$ at that level, and a union bound over the $N$ retained splits gives $P(E_1^{\mathrm c}\mid\mathcal G)\le\delta$,} while Lemma~\ref{lem:random-subsampling-plurality} gives $P(E_2^{\mathrm c}\mid\mathcal G)\le P(R=0\mid\mathcal G)+L\,\mathbb E\{e^{-R\Delta_\Pi^2/2}\mid\mathcal G\}$. On $E_1$, two retained splits in the same cell satisfy $|\widehat\theta_S-\widehat\theta_{S'}|\le2(\rho-\varepsilon)+2\varepsilon=2\rho$, so they are linked. Two retained splits in different cells satisfy $|\widehat\theta_S-\widehat\theta_{S'}|>2(\rho+\varepsilon)-2\varepsilon=2\rho$, so they are not linked. Hence the connected components are exactly the nonempty sets $\mathcal R_m\cap C_\ell$. On $E_2$, $\mathcal R_m\cap C_0$ is the unique largest component, and Algorithm~\ref{alg:probe-search} returns the median of its estimates. For every $S\in C_0$, $|\widehat\theta_S-\tau|\le|\widehat\theta_S-\theta_S|+|\theta_S-\tau|\le\varepsilon+b$\textcolor{blue}{, where $|\theta_S-\tau|\le b$ by Eq.~\eqref{eq:uniform-approximate-proxy-bias-bound}, because every $S\in C_0$ satisfies Assumptions~\ref{ass:consistency}, \ref{ass:joint-latent-exchangeability}, and~\ref{ass:uniform-stability}}. The median therefore lies in $[\tau-(b+\varepsilon),\tau+(b+\varepsilon)]$. A union bound over $E_1^{\mathrm c}$ and $E_2^{\mathrm c}$ proves Eq.~\eqref{eq:probe-search-error-probability}.
\end{proof}

\begin{proof}[Proof of Corollary~\ref{cor:budget-rule}]
Write $c\triangleq\Delta_\Pi^2/2$. First, $P(R=0)=\prod_{i=0}^{m-1}(T-N-i)/(T-i)\le(1-p)^m\le e^{-mp}$. Second, $R$ is the sum of the indicators $\mathbf 1\{S\in\widehat{\mathcal B}\}$ over the $m$ splits drawn without replacement. Because $x\mapsto e^{-cx}$ is continuous and convex, Thm.~4 of \citet{hoeffding1963probability} bounds $\mathbb E\{e^{-cR}\}$ by the corresponding expectation under sampling with replacement. Thus $\mathbb E\{e^{-cR}\}\le\{1-p(1-e^{-c})\}^m\le\exp\{-mp(1-e^{-c})\}$. Since $1-e^{-c}\le1$, adding the bound for $P(R=0)$ gives Eq.~\eqref{eq:sampling-failure-closed-form}. Its right side is at most $\alpha$ when $m\ge\log\{(L+1)/\alpha\}/\{p(1-e^{-c})\}$, and Thm.~\ref{thm:probe-search-error} gives $P(|\widehat\tau-\tau|\le b+\varepsilon)\ge1-\delta-\alpha$. Finally, $1-e^{-c}\ge c/(1+c)\ge\Delta_\Pi^2/3$ because $0<\Delta_\Pi\le1$. If $L=0$, only the empty-draw bound is needed, and $m\ge\log(1/\alpha)/p$ suffices.
\end{proof}

\paragraph{Median under majority.}
\begin{proof}[Proof of Corollary~\ref{cor:median-under-majority}]
On the target-band and uniform-evaluation-error event, every retained split in $C_0$ satisfies $|\widehat\theta_S-\tau|\le b+\varepsilon$. If $|\mathcal R_m\cap C_0|>R/2$, more than half of all retained estimates lie in $[\tau-(b+\varepsilon),\tau+(b+\varepsilon)]$. Both middle order statistics, and hence their midpoint when needed, lie in this interval. Therefore the median lies within $b+\varepsilon$ of $\tau$.
\end{proof}

\subsection{Finite-Strata MMD Specialization}
\label{app:finite-strata-mmd-specialization}

This subsection records the earlier finite-strata analysis as an independent specialization. It is not used by the continuous-representation results in the main text.

Fix a finite library $\mathcal F_j^{\mathrm{fs}}\subseteq\Phi_j$ with $N_j^{\mathrm{fs}}\triangleq|\mathcal F_j^{\mathrm{fs}}|$. Every $\phi\in\mathcal F_j^{\mathrm{fs}}$ is deterministic and takes values in a common finite set $\mathcal S_j$ of size $M_j$. Write $Z_j^\phi=\phi(V_j)$, let $\mathcal R_j(\phi)$ be its positive-mass strata, and set $p_{j,z}(\phi)=P(Z_j^\phi=z)$. Assume, uniformly over the library and $z\in\mathcal R_j(\phi)$,
\(P(Z_j^\phi=z)\ge p_{\min}>0, \eta\le P(A=1\mid Z_j^\phi=z)\le1-\eta\)
for some $\eta\in(0,1/2]$. Let $\kappa_j$ be a characteristic kernel on $(X,W_j)$ with reproducing kernel Hilbert space $\mathcal H_j$ and
\(\sup_v\sqrt{\kappa_j(v,v)}\le\rho_j<\infty.\)
For $z\in\mathcal R_j(\phi)$, define
\(\Delta_{j,z}^{\mathrm{MMD}}(\phi) \triangleq \operatorname{MMD}_{\kappa_j} \left\{ \textcolor{red}{\mathcal L}\textcolor{blue}{P}(X,W_j\mid A=1,Z_j^\phi=z), \textcolor{red}{\mathcal L}\textcolor{blue}{P}(X,W_j\mid A=0,Z_j^\phi=z) \right\},\); \(D_j^{\mathrm{MMD}}(\phi) \triangleq \sum_{z\in\mathcal R_j(\phi)} p_{j,z}(\phi)\Delta_{j,z}^{\mathrm{MMD}}(\phi).\)

For $i\in\mathcal I_{\mathrm D}$, let $Z_{j,i}^\phi=\phi(V_{j,i})$ and define
\(N_{a,z}^\phi \triangleq \sum_{i\in\mathcal I_{\mathrm D}} \mathbf 1\{A_i=a,Z_{j,i}^\phi=z\},\); \(N_z^\phi\triangleq N_{0,z}^\phi+N_{1,z}^\phi,\); \(\widehat p_{j,z}(\phi)\triangleq\frac{N_z^\phi}{n_{\mathrm D}},\); \(\widehat m_{a,j,z}(\phi) \triangleq \frac{1}{N_{a,z}^\phi} \sum_{i\in\mathcal I_{\mathrm D}} \mathbf 1\{A_i=a,Z_{j,i}^\phi=z\}\); \(\qquad{} \times\kappa_j\{(X_i,W_{j,i}),\cdot\}.\)
When both treatment-arm cells are nonempty, set
\(\widehat\Delta_{j,z}^{\mathrm{MMD}}(\phi) \triangleq \|\widehat m_{1,j,z}(\phi)-\widehat m_{0,j,z}(\phi)\|_{\mathcal H_j},\); \(\widehat D_j^{\mathrm{MMD}}(\phi) \triangleq \sum_{z:N_z^\phi>0} \widehat p_{j,z}(\phi)\widehat\Delta_{j,z}^{\mathrm{MMD}}(\phi).\)
If an occupied stratum contains only one treatment arm, set $\widehat D_j^{\mathrm{MMD}}(\phi)=+\infty$.

The MMD specialization uses its own structural stability condition. For every $\phi\in\mathcal F_j^{\mathrm{fs}}$, every positive-mass stratum $z$, every $a\in\{0,1\}$, and all $Q,Q'\in\mathcal Q_{j,z}$, define
\(\Delta^\mu_{a,j,z}(Q,Q') \triangleq \mathbb E_Q\{\mu_{a,j,z}\}-\mathbb E_{Q'}\{\mu_{a,j,z}\}.\)
Assume
\begin{align}
\left|\Delta^\mu_{a,j,z}(Q,Q')\right|
\le
\Gamma_j^{\mathrm{MMD}}(\phi)
\operatorname{MMD}_{\kappa_j}(\mathcal K_jQ,\mathcal K_jQ'),
\qquad
\sup_{\phi\in\mathcal F_j^{\mathrm{fs}}}
\Gamma_j^{\mathrm{MMD}}(\phi)
\le\overline\Gamma_j^{\mathrm{MMD}}<\infty.
\label{eq:finite-mmd-stability}
\end{align}

\begin{proposition}[Finite-strata MMD learning bound]
\label{prop:finite-strata-mmd-learning-bound}
Let $\delta_{\mathrm D}\in(0,1)$ and suppose
\begin{align}
n_{\mathrm D}
\ge
\frac{8}{p_{\min}\eta}
\log\left(\frac{8M_jN_j^{\mathrm{fs}}}{\delta_{\mathrm D}}\right).
\label{eq:finite-mmd-sample-threshold}
\end{align}
Define
\(e_{j,n_{\mathrm D}}^{\mathrm{MMD}}(\delta_{\mathrm D}) \triangleq \rho_j \left[ 2+\sqrt{2\log\left(\frac{4M_jN_j^{\mathrm{fs}}}{\delta_{\mathrm D}}\right)} \right] \sqrt{\frac{2}{n_{\mathrm D}p_{\min}\eta}},\); \(r_{j,n_{\mathrm D}}^{\mathrm{MMD}}(\delta_{\mathrm D}) \triangleq 2e_{j,n_{\mathrm D}}^{\mathrm{MMD}}(\delta_{\mathrm D}) +2\rho_j \sqrt{ \frac{ 2\left[M_j\log2+\log\{4N_j^{\mathrm{fs}}/\delta_{\mathrm D}\}\right] }{n_{\mathrm D}} }.\)
Then, with probability at least $1-\delta_{\mathrm D}$,
\begin{align}
\sup_{\phi\in\mathcal F_j^{\mathrm{fs}}}
\left|
\widehat D_j^{\mathrm{MMD}}(\phi)-D_j^{\mathrm{MMD}}(\phi)
\right|
\le
r_{j,n_{\mathrm D}}^{\mathrm{MMD}}(\delta_{\mathrm D}).
\label{eq:finite-mmd-uniform-deviation}
\end{align}
If $\widehat\phi_j^{\mathrm{fs}}$ obeys
\(\widehat D_j^{\mathrm{MMD}}(\widehat\phi_j^{\mathrm{fs}}) \le \inf_{\phi\in\mathcal F_j^{\mathrm{fs}}} \widehat D_j^{\mathrm{MMD}}(\phi) +\varepsilon_{\mathrm{opt},j}^{\mathrm{fs}},\)
then, on the same event,
\begin{align}
D_j^{\mathrm{MMD}}(\widehat\phi_j^{\mathrm{fs}})
&\le
\inf_{\phi\in\mathcal F_j^{\mathrm{fs}}}
D_j^{\mathrm{MMD}}(\phi)
+2r_{j,n_{\mathrm D}}^{\mathrm{MMD}}(\delta_{\mathrm D})
+\varepsilon_{\mathrm{opt},j}^{\mathrm{fs}},
\label{eq:finite-mmd-oracle}\\
|\tau_{\widehat Z_j^{\mathrm{fs}}}-\tau|
&\le
\overline\Gamma_j^{\mathrm{MMD}}
\left[
\inf_{\phi\in\mathcal F_j^{\mathrm{fs}}}
D_j^{\mathrm{MMD}}(\phi)
+2r_{j,n_{\mathrm D}}^{\mathrm{MMD}}(\delta_{\mathrm D})
+\varepsilon_{\mathrm{opt},j}^{\mathrm{fs}}
\right],
\label{eq:finite-mmd-causal-oracle}
\end{align}
where $\widehat Z_j^{\mathrm{fs}}=\widehat\phi_j^{\mathrm{fs}}(V_j)$, provided the primitive structural conditions and overlap used in the exact bias decomposition hold for every candidate.
\end{proposition}

\begin{proof}
Eq.~\eqref{eq:finite-mmd-sample-threshold} and a multiplicative Chernoff bound imply that every candidate treatment-by-stratum cell contains at least $n_{\mathrm D}p_{\min}\eta/2$ observations except on an event of probability at most $\delta_{\mathrm D}/4$. Conditional on the cell assignments, bounded Hilbert-space mean concentration and a union bound over at most $2M_jN_j^{\mathrm{fs}}$ cells give
\(\left| \widehat\Delta_{j,z}^{\mathrm{MMD}}(\phi) -\Delta_{j,z}^{\mathrm{MMD}}(\phi) \right| \le 2e_{j,n_{\mathrm D}}^{\mathrm{MMD}}(\delta_{\mathrm D})\)
simultaneously, with failure probability at most $\delta_{\mathrm D}/2$. A multinomial $L_1$ concentration bound gives
\(\sum_z|\widehat p_{j,z}(\phi)-p_{j,z}(\phi)| \le \sqrt{ \frac{ 2\left[M_j\log2+\log\{4N_j^{\mathrm{fs}}/\delta_{\mathrm D}\}\right] }{n_{\mathrm D}} }\)
simultaneously over candidates, with failure probability at most $\delta_{\mathrm D}/4$. On the intersection of these events, $0\le\Delta_{j,z}^{\mathrm{MMD}}(\phi)\le2\rho_j$ and the triangle inequality give Eq.~\eqref{eq:finite-mmd-uniform-deviation}.

The empirical-risk comparison yields
\(D_j^{\mathrm{MMD}}(\widehat\phi_j^{\mathrm{fs}}) \le \widehat D_j^{\mathrm{MMD}}(\widehat\phi_j^{\mathrm{fs}}) +r_{j,n_{\mathrm D}}^{\mathrm{MMD}}(\delta_{\mathrm D}) \le \inf_{\phi\in\mathcal F_j^{\mathrm{fs}}} D_j^{\mathrm{MMD}}(\phi) +2r_{j,n_{\mathrm D}}^{\mathrm{MMD}}(\delta_{\mathrm D}) +\varepsilon_{\mathrm{opt},j}^{\mathrm{fs}}\),
which proves Eq.~\eqref{eq:finite-mmd-oracle}. The exact bias decomposition and the MMD stability condition in Eq.~\eqref{eq:finite-mmd-stability} give
\(|\tau_{Z_j^\phi}-\tau| \le \Gamma_j^{\mathrm{MMD}}(\phi)D_j^{\mathrm{MMD}}(\phi).\)
Apply this inequality to $\widehat\phi_j^{\mathrm{fs}}$ and substitute Eq.~\eqref{eq:finite-mmd-oracle} to obtain Eq.~\eqref{eq:finite-mmd-causal-oracle}.
\end{proof}
